\documentclass{article}

\usepackage{../handouts}

\title{CSCI 432 Handout 04: Recurrence Relations}
\author{Name: \rule{3in}{0.15mm} \\~\\
    Collaborators: \rule{3in}{0.15mm} }
\date{27 January 2026}

\begin{document}
\maketitle

\section*{Recall a Couple Common Summations}

There are a couple summations that we see over and over again in algorithms.
Let's see which you remember.  Please fill out the following table, where all
functions are of the form $f_i \colon \N \to \R$, and $c \in \R$ is a constant.

\begin{table}[h!]
    \begin{tabular}{p{0.75in} p{1.25in} p{2in} p{2in}}
        \toprule
        Name & Summation    & Closed Form & Asymptotic Form \\
        \midrule
        Arithmetic Series   & $f_1(n) = \sum_{i=0}^n c$ \vspace{.9in} & & \\
        Constant Summation  & $f_2(n) = \sum_{i=0}^n i$ \vspace{.9in} & & \\
        Sum of Squares  & $f_3(n) = \sum_{i=1}^n i^2$ \vspace{.9in} & & \\
        Finite Geom.\ Series & $f_4(n) = \sum_{i=0}^n c r^i$ \vspace{.9in} &
            $f_4(n) = \begin{cases}
                c\frac{1-r^{n+1}}{1-r} & r \neq 1 \\
                ? & r = 1
                    \end{cases}$                         & \\
        Harmonic Series  & $f_5(n) = \sum_{i=1}^n \frac{1}{i}$ \vspace{.9in} &
            (no nice closed form) & \\
        \bottomrule
    \end{tabular}
\end{table}

\FloatBarrier
\section*{Summations and Recursions}

Now, let's write each of the summations above as a recursion by `plucking off'
the first term.  For example,

$$ f_1(n) = \begin{cases}
                c & n=1 \\
                c + f_1(n-1) & otw
            \end{cases}
$$

\begin{enumerate}
    \item $f_2(n) = $
        \vspace{1in}
    \item $f_3(n) = $
        \vspace{1in}
    \item $f_4(n) = $
        \vspace{1in}
    \item $f_5(n) = $
        \vspace{1in}
\end{enumerate}

Being able to do this is helpful for:
\begin{itemize}
    \item In the inductive step of a proof by induction.
    \item To walk through algorithms.
\end{itemize}


\section*{Recursion}

For each of the recursions below, name an algorithm that uses that recursion, and
give the asymptotic form of the recurrence relation.

\begin{enumerate}
    \item $T(n) = 2 T(n/2) + \Theta(n)$\\
        \answer
    \item $T(n) = T(n/2) + \Theta(1)$\\
        \answer
    \item $T(n) = T(n-1) + \Theta(1)$\\
        \answer
    \item $T(n) = T(n-1) + \Theta(n)$\\
        \answer
    \item $T(n) = 2T(n-1) + \Theta(1)$\\
        \answer
    \item $T(n) = T(n-1) + T(n-2) + \Theta(1)$\\
        \answer
\end{enumerate}

We often dived the problem in half.  What happens in the following two
recursions if we divide the problem into thirds?
\begin{enumerate}
    \item $T(n) = 3 T(n/3) + \Theta(n)$\\
        \answer
    \item $T(n) = T(n/3) + \Theta(1)$\\
        \answer
\end{enumerate}

\end{document}
